\centerline{\bf \Huge ДЗ. Фальшивые монетки}
\centerline{\bf \Huge }

\begin{flushright}
        \scriptsize{\textbf{Вежливость, как игорная марка, есть открыто признанная фальшивая монета}}
\end{flushright}

\problem{\textbf{1}} Имеется 5 монет, внешне не отличающихся друг от друга. За какое наименьшее количество взвешиваний на чашечных весах можно гарантированно найти одну фальшивую монетку (ФМ), если известно, что она весит тяжелее настоящих.

\problem{\textbf{2}} Имеется 6 монет, внешне не отличающихся друг от друга. За какое наименьшее количество взвешиваний на чашечных весах можно гарантированно найти одну фальшивую монетку, если известно, что она может весить больше, а может и меньше настоящих.

\problem{\textbf{3}} Имеется 3 белых и 6 жёлтых монет. За какое наименьшее количество взвешиваний на чашечных весах можно гарантированно найти одну фальшивую монетку, если известно, что ФМ белая весит меньше настоящих белых, а ФМ жёлтая весит больше настоящих жёлтых. Белые настоящие и жёлтые настоящие имеют разный вес.

\problem{\textbf{4}} Имеется 3 белых, 4 жёлтых и 5 зелёных монет. За какое наименьшее количество взвешиваний на чашечных весах можно гарантированно найти одну фальшивую монетку, если известно, что ФМ белая весит меньше настоящих белых, ФМ жёлтая весит больше настоящих жёлтых, ФМ зелёная весит меньше настоящих зелёных. Белые настоящие, жёлтые настоящие и зелёные настоящие имеют разный вес.

